\documentclass[german, @CLASSOPTIONS@]{tumexam}

\usepackage{tummath}

\title{tumexam: Eine Prüfungsvorlage im TUM Corporate Design}
\subtitle{Abschlussklausur im Wahlmodul}

\author[orcid=0000-0002-5164-432X]{Dr.-Ing.\ Michael Haider}

\date{Freitag, den 29. April 2022\\
  13:15 -- 14:45 Uhr}

\studentblock{%
  \renewcommand{\arraystretch}{1.25}
  \begin{tabular}{l p{10cm}}
  Nachname:     & \hrule \\
  Vorname:      & \hrule \\
  Matrikel-Nr.: & \hrule \\
  Hörsaal:      & \hrule \\
  Platz-Nr.:    & \hrule
  \end{tabular}%
}

\disclaimer{%
  \begin{itemize}
  \item Dieses Aufgabenheft hat \textbf{\pageref{LastPage} Seiten}.
  \item Die Gesamtzahl der Punkte beträgt \numpoints.
  \item Die Bearbeitungszeit umfasst \textbf{90 Minuten}.
  \item Es sind \textbf{keine} Hilfsmittel erlaubt.
  \item Schriftliche Unterlagen und elektronische Geräte (Computer, Taschenrechner,
  Mobiltelefone usw.) sind \textbf{nicht} zugelassen.
  \item Mit * gekennzeichnete Aufgaben sind ohne Kenntnis der Ergebnisse der
  vorhergehenden Teilaufgaben lösbar.
  \item Benutzen Sie die vorgesehenen Boxen für Ihre Antworten. Falls Sie mehr Platz
  benötigen, benutzen Sie bitte die \textbf{jeweilige Rückseite} oder das
  \textbf{Zusatzblatt} am Ende des Aufgabenheftes (bitte Aufgaben-Nummer
  kennzeichnen!).
  \item \textbf{Gewertet werden nur Ergebnisse, bei denen der Lösungsweg
    erkennbar ist!}
  \end{itemize}%
}

\begin{document}

\maketitle

\clearpage

\begin{questions}
  %
  \titledquestion{Kurzfragen}[48]
  %
  \begin{parts}
    %
    \independentpart
    Wie erstellt man eine unabhängige Frage mit tumexam?
    %
    \begin{solutionbox}{3.25cm}
      %
      Eine unabhängige Frage kann mithilfe von \texttt{\textbackslash{independentpart}}, in einer \texttt{parts} Umgebung, in einer \texttt{questions} Umgebung erzeugt werden. Für weitere Details verweisen wir auf die Dokumentation der \LaTeX{} exam Klasse.
      %
    \end{solutionbox}
    %
    \dependentpart
    Wie erstellt man eine abhängige Frage mit tumexam?
    %
    \begin{solutionbox}{3.25cm}
      %
      Eine unabhängige Frage kann mithilfe von \texttt{\textbackslash{dependentpart}}, in einer \texttt{parts} Umgebung, in einer \texttt{questions} Umgebung erzeugt werden. Für weitere Details verweisen wir auf die Dokumentation der \LaTeX{} exam Klasse.
      %
    \end{solutionbox}
    %
  \end{parts}
  %
  \clearpage
  %
  \titledquestion{Ein größeres Problem}[42]
  %
  Angenommen wir möchten den Sinn des Lebens, also die Antwort auf das Leben, das Universum und überhaupt alles bestimmen. Beachten Sie dazu Abbildung~\ref{fig-graph-1}.
  %
  \begin{figure}[h]
    %
    \centering
    \begin{tikzpicture}
      %
      \draw (0,0) circle (2cm);
      %
    \end{tikzpicture}
    %
    \caption{Abbildung zum Sinn des Lebens.}
    %
    \label{fig-graph-1}
    %
  \end{figure}
  %
  \begin{parts}
    %
    \independentpart
    Was ist die Antwort auf das Leben, das Universum und überhaupt alles?
    %
    \begin{solutionbox}{4.5cm} % (1p)
      %
      Die Antwort auf das Leben, das Universum und überhaupt alles ist
      %
      \begin{equation*}
        42\, .
      \end{equation*}
      %
    \end{solutionbox}
    %
  \end{parts}
  %
  \clearpage
  %
  \titledquestion{Zusätzliche Fragen}[2]
  %
  \begin{parts}
    %
    \independentpart
    Wie erstellt man eine unabhängige Frage mit tumexam?
    %
    \begin{solutionbox}{3.25cm}
      %
      Eine unabhängige Frage kann mithilfe von \texttt{\textbackslash{independentpart}}, in einer \texttt{parts} Umgebung, in einer \texttt{questions} Umgebung erzeugt werden. Für weitere Details verweisen wir auf die Dokumentation der \LaTeX{} exam Klasse.
      %
    \end{solutionbox}
    %
    \dependentpart
    Wie erstellt man eine abhängige Frage mit tumexam?
    %
    \begin{solutionbox}{3.25cm}
      %
      Eine unabhängige Frage kann mithilfe von \texttt{\textbackslash{dependentpart}}, in einer \texttt{parts} Umgebung, in einer \texttt{questions} Umgebung erzeugt werden. Für weitere Details verweisen wir auf die Dokumentation der \LaTeX{} exam Klasse.
      %
    \end{solutionbox}
    %
  \end{parts}
  %
  \clearpage
  %
  \titledquestion{Zusätzliche Fragen}[2]
  %
  \begin{parts}
    %
    \independentpart
    Wie erstellt man eine unabhängige Frage mit tumexam?
    %
    \begin{solutionbox}{3.25cm}
      %
      Eine unabhängige Frage kann mithilfe von \texttt{\textbackslash{independentpart}}, in einer \texttt{parts} Umgebung, in einer \texttt{questions} Umgebung erzeugt werden. Für weitere Details verweisen wir auf die Dokumentation der \LaTeX{} exam Klasse.
      %
    \end{solutionbox}
    %
    \dependentpart
    Wie erstellt man eine abhängige Frage mit tumexam?
    %
    \begin{solutionbox}{3.25cm}
      %
      Eine unabhängige Frage kann mithilfe von \texttt{\textbackslash{dependentpart}}, in einer \texttt{parts} Umgebung, in einer \texttt{questions} Umgebung erzeugt werden. Für weitere Details verweisen wir auf die Dokumentation der \LaTeX{} exam Klasse.
      %
    \end{solutionbox}
    %
  \end{parts}
  %
  \clearpage
  %
  \titledquestion{Zusätzliche Fragen}[2]
  %
  \begin{parts}
    %
    \independentpart
    Wie erstellt man eine unabhängige Frage mit tumexam?
    %
    \begin{solutionbox}{3.25cm}
      %
      Eine unabhängige Frage kann mithilfe von \texttt{\textbackslash{independentpart}}, in einer \texttt{parts} Umgebung, in einer \texttt{questions} Umgebung erzeugt werden. Für weitere Details verweisen wir auf die Dokumentation der \LaTeX{} exam Klasse.
      %
    \end{solutionbox}
    %
    \dependentpart
    Wie erstellt man eine abhängige Frage mit tumexam?
    %
    \begin{solutionbox}{3.25cm}
      %
      Eine unabhängige Frage kann mithilfe von \texttt{\textbackslash{dependentpart}}, in einer \texttt{parts} Umgebung, in einer \texttt{questions} Umgebung erzeugt werden. Für weitere Details verweisen wir auf die Dokumentation der \LaTeX{} exam Klasse.
      %
    \end{solutionbox}
    %
  \end{parts}
  %
  \clearpage
  %
  \titledquestion{Zusätzliche Fragen}[2]
  %
  \begin{parts}
    %
    \independentpart
    Wie erstellt man eine unabhängige Frage mit tumexam?
    %
    \begin{solutionbox}{3.25cm}
      %
      Eine unabhängige Frage kann mithilfe von \texttt{\textbackslash{independentpart}}, in einer \texttt{parts} Umgebung, in einer \texttt{questions} Umgebung erzeugt werden. Für weitere Details verweisen wir auf die Dokumentation der \LaTeX{} exam Klasse.
      %
    \end{solutionbox}
    %
    \dependentpart
    Wie erstellt man eine abhängige Frage mit tumexam?
    %
    \begin{solutionbox}{3.25cm}
      %
      Eine unabhängige Frage kann mithilfe von \texttt{\textbackslash{dependentpart}}, in einer \texttt{parts} Umgebung, in einer \texttt{questions} Umgebung erzeugt werden. Für weitere Details verweisen wir auf die Dokumentation der \LaTeX{} exam Klasse.
      %
    \end{solutionbox}
    %
  \end{parts}
  %
  \clearpage
  %
  \titledquestion{Zusätzliche Fragen}[2]
  %
  \begin{parts}
    %
    \independentpart
    Wie erstellt man eine unabhängige Frage mit tumexam?
    %
    \begin{solutionbox}{3.25cm}
      %
      Eine unabhängige Frage kann mithilfe von \texttt{\textbackslash{independentpart}}, in einer \texttt{parts} Umgebung, in einer \texttt{questions} Umgebung erzeugt werden. Für weitere Details verweisen wir auf die Dokumentation der \LaTeX{} exam Klasse.
      %
    \end{solutionbox}
    %
    \dependentpart
    Wie erstellt man eine abhängige Frage mit tumexam?
    %
    \begin{solutionbox}{3.25cm}
      %
      Eine unabhängige Frage kann mithilfe von \texttt{\textbackslash{dependentpart}}, in einer \texttt{parts} Umgebung, in einer \texttt{questions} Umgebung erzeugt werden. Für weitere Details verweisen wir auf die Dokumentation der \LaTeX{} exam Klasse.
      %
    \end{solutionbox}
    %
  \end{parts}
  %
  \clearpage
  %
  \titledquestion{Zusätzliche Fragen}[2]
  %
  \begin{parts}
    %
    \independentpart
    Wie erstellt man eine unabhängige Frage mit tumexam?
    %
    \begin{solutionbox}{3.25cm}
      %
      Eine unabhängige Frage kann mithilfe von \texttt{\textbackslash{independentpart}}, in einer \texttt{parts} Umgebung, in einer \texttt{questions} Umgebung erzeugt werden. Für weitere Details verweisen wir auf die Dokumentation der \LaTeX{} exam Klasse.
      %
    \end{solutionbox}
    %
    \dependentpart
    Wie erstellt man eine abhängige Frage mit tumexam?
    %
    \begin{solutionbox}{3.25cm}
      %
      Eine unabhängige Frage kann mithilfe von \texttt{\textbackslash{dependentpart}}, in einer \texttt{parts} Umgebung, in einer \texttt{questions} Umgebung erzeugt werden. Für weitere Details verweisen wir auf die Dokumentation der \LaTeX{} exam Klasse.
      %
    \end{solutionbox}
    %
  \end{parts}
  %
  \clearpage
  %
  \titledquestion{Zusätzliche Fragen}[2]
  %
  \begin{parts}
    %
    \independentpart
    Wie erstellt man eine unabhängige Frage mit tumexam?
    %
    \begin{solutionbox}{3.25cm}
      %
      Eine unabhängige Frage kann mithilfe von \texttt{\textbackslash{independentpart}}, in einer \texttt{parts} Umgebung, in einer \texttt{questions} Umgebung erzeugt werden. Für weitere Details verweisen wir auf die Dokumentation der \LaTeX{} exam Klasse.
      %
    \end{solutionbox}
    %
    \dependentpart
    Wie erstellt man eine abhängige Frage mit tumexam?
    %
    \begin{solutionbox}{3.25cm}
      %
      Eine unabhängige Frage kann mithilfe von \texttt{\textbackslash{dependentpart}}, in einer \texttt{parts} Umgebung, in einer \texttt{questions} Umgebung erzeugt werden. Für weitere Details verweisen wir auf die Dokumentation der \LaTeX{} exam Klasse.
      %
    \end{solutionbox}
    %
  \end{parts}
  %
  \clearpage
  %
  \titledquestion{Zusätzliche Fragen}[2]
  %
  \begin{parts}
    %
    \independentpart
    Wie erstellt man eine unabhängige Frage mit tumexam?
    %
    \begin{solutionbox}{3.25cm}
      %
      Eine unabhängige Frage kann mithilfe von \texttt{\textbackslash{independentpart}}, in einer \texttt{parts} Umgebung, in einer \texttt{questions} Umgebung erzeugt werden. Für weitere Details verweisen wir auf die Dokumentation der \LaTeX{} exam Klasse.
      %
    \end{solutionbox}
    %
    \dependentpart
    Wie erstellt man eine abhängige Frage mit tumexam?
    %
    \begin{solutionbox}{3.25cm}
      %
      Eine unabhängige Frage kann mithilfe von \texttt{\textbackslash{dependentpart}}, in einer \texttt{parts} Umgebung, in einer \texttt{questions} Umgebung erzeugt werden. Für weitere Details verweisen wir auf die Dokumentation der \LaTeX{} exam Klasse.
      %
    \end{solutionbox}
    %
  \end{parts}
  %
  \clearpage
  %
  \titledquestion{Zusätzliche Fragen}[2]
  %
  \begin{parts}
    %
    \independentpart
    Wie erstellt man eine unabhängige Frage mit tumexam?
    %
    \begin{solutionbox}{3.25cm}
      %
      Eine unabhängige Frage kann mithilfe von \texttt{\textbackslash{independentpart}}, in einer \texttt{parts} Umgebung, in einer \texttt{questions} Umgebung erzeugt werden. Für weitere Details verweisen wir auf die Dokumentation der \LaTeX{} exam Klasse.
      %
    \end{solutionbox}
    %
    \dependentpart
    Wie erstellt man eine abhängige Frage mit tumexam?
    %
    \begin{solutionbox}{3.25cm}
      %
      Eine unabhängige Frage kann mithilfe von \texttt{\textbackslash{dependentpart}}, in einer \texttt{parts} Umgebung, in einer \texttt{questions} Umgebung erzeugt werden. Für weitere Details verweisen wir auf die Dokumentation der \LaTeX{} exam Klasse.
      %
    \end{solutionbox}
    %
  \end{parts}
  %
  \clearpage
  %
  \titledquestion{Zusätzliche Fragen}[2]
  %
  \begin{parts}
    %
    \independentpart
    Wie erstellt man eine unabhängige Frage mit tumexam?
    %
    \begin{solutionbox}{3.25cm}
      %
      Eine unabhängige Frage kann mithilfe von \texttt{\textbackslash{independentpart}}, in einer \texttt{parts} Umgebung, in einer \texttt{questions} Umgebung erzeugt werden. Für weitere Details verweisen wir auf die Dokumentation der \LaTeX{} exam Klasse.
      %
    \end{solutionbox}
    %
    \dependentpart
    Wie erstellt man eine abhängige Frage mit tumexam?
    %
    \begin{solutionbox}{3.25cm}
      %
      Eine unabhängige Frage kann mithilfe von \texttt{\textbackslash{dependentpart}}, in einer \texttt{parts} Umgebung, in einer \texttt{questions} Umgebung erzeugt werden. Für weitere Details verweisen wir auf die Dokumentation der \LaTeX{} exam Klasse.
      %
    \end{solutionbox}
    %
  \end{parts}
  %
  \clearpage
  %
  \titledquestion{Zusätzliche Fragen}[2]
  %
  \begin{parts}
    %
    \independentpart
    Wie erstellt man eine unabhängige Frage mit tumexam?
    %
    \begin{solutionbox}{3.25cm}
      %
      Eine unabhängige Frage kann mithilfe von \texttt{\textbackslash{independentpart}}, in einer \texttt{parts} Umgebung, in einer \texttt{questions} Umgebung erzeugt werden. Für weitere Details verweisen wir auf die Dokumentation der \LaTeX{} exam Klasse.
      %
    \end{solutionbox}
    %
    \dependentpart
    Wie erstellt man eine abhängige Frage mit tumexam?
    %
    \begin{solutionbox}{3.25cm}
      %
      Eine unabhängige Frage kann mithilfe von \texttt{\textbackslash{dependentpart}}, in einer \texttt{parts} Umgebung, in einer \texttt{questions} Umgebung erzeugt werden. Für weitere Details verweisen wir auf die Dokumentation der \LaTeX{} exam Klasse.
      %
    \end{solutionbox}
    %
  \end{parts}
  %
  \clearpage
  %
  \titledquestion{Zusätzliche Fragen}[2]
  %
  \begin{parts}
    %
    \independentpart
    Wie erstellt man eine unabhängige Frage mit tumexam?
    %
    \begin{solutionbox}{3.25cm}
      %
      Eine unabhängige Frage kann mithilfe von \texttt{\textbackslash{independentpart}}, in einer \texttt{parts} Umgebung, in einer \texttt{questions} Umgebung erzeugt werden. Für weitere Details verweisen wir auf die Dokumentation der \LaTeX{} exam Klasse.
      %
    \end{solutionbox}
    %
    \dependentpart
    Wie erstellt man eine abhängige Frage mit tumexam?
    %
    \begin{solutionbox}{3.25cm}
      %
      Eine unabhängige Frage kann mithilfe von \texttt{\textbackslash{dependentpart}}, in einer \texttt{parts} Umgebung, in einer \texttt{questions} Umgebung erzeugt werden. Für weitere Details verweisen wir auf die Dokumentation der \LaTeX{} exam Klasse.
      %
    \end{solutionbox}
    %
  \end{parts}
  %
  \clearpage
  %
  \titledquestion{Zusätzliche Fragen}[2]
  %
  \begin{parts}
    %
    \independentpart
    Wie erstellt man eine unabhängige Frage mit tumexam?
    %
    \begin{solutionbox}{3.25cm}
      %
      Eine unabhängige Frage kann mithilfe von \texttt{\textbackslash{independentpart}}, in einer \texttt{parts} Umgebung, in einer \texttt{questions} Umgebung erzeugt werden. Für weitere Details verweisen wir auf die Dokumentation der \LaTeX{} exam Klasse.
      %
    \end{solutionbox}
    %
    \dependentpart
    Wie erstellt man eine abhängige Frage mit tumexam?
    %
    \begin{solutionbox}{3.25cm}
      %
      Eine unabhängige Frage kann mithilfe von \texttt{\textbackslash{dependentpart}}, in einer \texttt{parts} Umgebung, in einer \texttt{questions} Umgebung erzeugt werden. Für weitere Details verweisen wir auf die Dokumentation der \LaTeX{} exam Klasse.
      %
    \end{solutionbox}
    %
  \end{parts}
  %
  \clearpage
  %
  \makenonemptybox{24.5cm}{\centering{(Zusatzblatt)}}
  %
\end{questions}

\end{document}
